% ============================================================
% §3 A Two-Level Framework for Noise Trajectory Search
% Content from method_framework.tex, restructured as independent section
% ============================================================
\section{A Two-Level Framework for Noise Trajectory Search}
\label{sec:framework}

We cast noise trajectory search as a two-level optimization problem.
A general transition abstraction covers all diffusion-family models;
the local search operator and global scheduler serve as its two components.
We conclude by showing how existing methods arise as special cases.

% -----------------------------------------------------------
\subsection{A General Transition Abstraction}
\label{sec:transition}

Every diffusion-family model generates a sample by traversing
$T$ discrete timesteps $t_T > t_{T-1} > \dots > t_1$.
At each step the sampler applies a transition
\begin{equation}
  x_{t-1} = F_\theta(x_t, t, c, \varepsilon_t),
  \label{eq:transition}
\end{equation}
where $x_t$ is the latent at timestep $t$, $c$ is the conditioning signal
(e.g., a class label or text prompt), and $\varepsilon_t$ is a noise variable
injecting stochasticity.
The final step yields $x_0$, which the verifier $v(x_0, c)$ then scores.

\paragraph{Applicability across model families.}
The transition in \cref{eq:transition} is deliberately model-agnostic.
For \emph{stochastic samplers} such as DDPM and SDE-based
models~\citep{ho2020denoising,song2021scorebased}, $\varepsilon_t$ is the
variance noise already present in sampling; searching over
$\varepsilon_t$ therefore extends the base sampler naturally.
For \emph{deterministic samplers} such as flow-matching ODE
solvers~\citep{lipman2023flow,liu2023flow}, the transition is nominally
noise-free.
Yet stochasticity can be injected at selected timesteps, by converting the ODE
into a family of SDEs sharing the same marginal
distributions~\citep{dockhorn2024stochastic}, or by adding controlled
perturbations for exploration~\citep{kim2025rbf,black2023training,fan2024reinforcement}.
Either way, sampling takes the form of
\cref{eq:transition}, so our framework applies to any diffusion-family
model.

% -----------------------------------------------------------
\subsection{The Two-Level View}
\label{sec:two-level}

Noise trajectory search modifies the noise
variables $\{\varepsilon_t\}_{t=1}^T$ along the latent trajectory
$\{x_t\}_{t=1}^T$ under a fixed test-time compute budget to
maximize the verifier score.
This problem decomposes naturally into two levels
(\cref{fig:architecture}):

\begin{figure}[t]
\centering
\includegraphics[width=\linewidth]{figures/fig_architecture.pdf}
\caption{Two-level view of noise trajectory search.
The \emph{global scheduler}~$\mathcal{G}$ (top) receives budget~$B$
and distributes per-step budgets $\{K_t\}$ across the denoising trajectory.
At each timestep, the \emph{local search operator}~$\mathcal{L}_t$ (bottom)
evaluates $K_t$ noise candidates and selects the best one.
Bars above each step depict the non-uniform allocation:
timesteps where noise perturbations matter more receive larger budgets.}
\label{fig:architecture}
\end{figure}

\paragraph{Local search operator.}
At timestep $t$, given the current latent $x_t$, conditioning $c$,
verifier $v$, and base sampler $F_\theta$, the \emph{local} search operator
\[
  \mathcal{L}_t : (x_t, c, v, F_\theta, \text{budget at } t)
  \longrightarrow (x_{t-1}, \varepsilon_t)
\]
samples multiple noise candidates $\varepsilon_t$ and selects the one that
maximizes the verifier score.
To evaluate candidate $\varepsilon_t^{(k)}$, we apply the
base sampler, obtaining the next-step latent and its clean-output
prediction in a single network pass:
\begin{equation}
  x_{t-1}^{(k)} = F_\theta(x_t,\, t,\, c,\, \varepsilon_t^{(k)}),
  \qquad
  \hat{x}_0^{(k)} = D_\theta(x_{t-1}^{(k)},\, t\!-\!1,\, c),
  \label{eq:x0_prediction}
\end{equation}
where $D_\theta$ denotes Tweedie's denoising estimate~\citep{song2021scorebased},
i.e., an intermediate network output that maps any noisy latent to a
clean-data estimate.
The verifier then scores this prediction:
\begin{equation}
  s^{(k)} = v\!\bigl(\hat{x}_0^{(k)},\, c\bigr).
  \label{eq:score}
\end{equation}
In one local search round, $K$ candidate noises
$\{\varepsilon_t^{(k)}\}_{k=1}^K$ are sampled and scored via
\cref{eq:x0_prediction,eq:score}; the winner
$k^* = \arg\max_k s^{(k)}$ is kept.
Each candidate costs one NFE, so $K$ candidates consume $K$ NFE.
Many existing search procedures (zero-order perturbation,
$\epsilon$-greedy exploration, particle-based
refinement) instantiate this operator; we treat
$\mathcal{L}_t$ abstractly as a black-box mapping from $x_t$ to $x_{t-1}$
given a per-step budget.

\paragraph{Global scheduler over timesteps.}
The local operator $\mathcal{L}_t$ decides \emph{how} to search at a given
step, but not \emph{where along the trajectory} to invest computation.
A global scheduler fills this role:
\[
  \mathcal{G} :
  \bigl(x_T, c, v, F_\theta, B, \{\mathcal{L}_t\}_{t=1}^T\bigr)
  \longrightarrow \bigl(x_0, \{x_t\}_{t=1}^T\bigr),
\]
which, starting from $x_T \sim \mathcal{N}(0, I)$, alternates
denoising steps with local search under budget $B$ and
returns both the final sample $x_0$ and the full latent trajectory
$\{x_t\}_{t=1}^T$.

Concretely, $\mathcal{G}$ walks timesteps from $T$ to $1$.
At each step $t$ it either
(i)~invokes $\mathcal{L}_t$ one or more times, spending
part of the remaining budget, or
(ii)~skips search and applies a single base step
$x_{t-1} = F_\theta(x_t, t, c, \varepsilon_t)$.
This decision must be made \emph{online}: at step $t$
we cannot predict how effective search will be at later timesteps, yet
computation spent now is lost forever.
Our goal is to design scheduling policies that decide, under a fixed NFE
budget, \emph{on which timesteps to search and for how long}.

% -----------------------------------------------------------
\subsection{Existing Methods as Special Cases}
\label{sec:special-cases}

This two-level view (local operator $\mathcal{L}_t$ plus global scheduler
$\mathcal{G}$) unifies many prior
inference-time search strategies:
\begin{itemize}
  \item \textbf{Plain diffusion sampling.}
  If $\mathcal{L}_t$ reduces to a single denoising step with no extra budget,
  $\mathcal{G}$ simply applies $F_\theta$ once per
  timestep, recovering standard sampling.
  \item \textbf{Best-of-$N$ / random search.}
  Drawing multiple i.i.d.\ initial noises, propagating each to $x_0$ without
  per-step modification, and selecting the best sample amounts to a
  scheduler that spends all compute on full trajectories, performing only
  a single global selection; no local search occurs.
  \item \textbf{Noise-trajectory and tree-style search.}
  Methods that iteratively refine a noise pivot or expand/prune partial
  trajectories choose specific
  forms of $\mathcal{L}_t$ paired with hard-coded $\mathcal{G}$ policies that
  fix which timesteps receive search and how budget splits across time.
\end{itemize}
The remainder of the paper focuses on a particular, practically
important family of methods, for which we develop a new global scheduler.
A more general MDP formulation that subsumes the above and accommodates richer
action spaces (multiple verifier types, backtracking) is presented in
\cref{sec:mdp}.
